\chapter{कार्य और यांत्रिक ऊर्जा}\label{ch:g12:work-and-energy}

सलाद के कटोरे में एक कंचा उड़ेल दीजिए: वह गोता लगाता है, तल पार कर जाता है,
अपनी शुरुआती ऊँचाई तक चढ़ता है, ठिठकता है, और लौट पड़ता है --- और हर पलटाव एक
वक्र (कटोरे की काट) और एक क्षैतिज रेखा (कंचे की ऊर्जा) से पहले ही बताया
जा सकता है, बिना कोई गति-समीकरण हल किए। पिछले वर्ष सीधी रेखाओं पर बलों
का दाम लगा था (\cref{ch:g11:work-of-force}) और गतिज तथा \omterm{def:g12:work-and-energy:ep}{स्थितिज ऊर्जा} का तराज़ू बना था
(\cref{ch:g11:mechanical-energy}); इस वर्ष वही बहीखाता हर जगह पहुँच जाता है: मुड़े हुए पथ, स्प्रिंगें,
और एक नज़र में पूरे भू-दृश्य।

\section{किसी भी पथ पर कार्य}

\begin{definition}[पथ पर कार्य]\label{def:g12:work-and-energy:path}
मान लीजिए $A$ से $B$ तक किसी भी पथ पर चलते पिंड पर बल $\vect F$ लगता
है: पथ को इतने छोटे विस्थापनों $\vect{\dd\ell}$ में काट दीजिए कि हर एक सीधा हो और उस पर
$\vect F$ मुश्किल से बदले। हर क़दम छोटा-सा कार्य $\scal{F}{\dd\ell} = F\,\dd\ell\cos\theta$ कमाता है (\cref{ch:g11:work-of-force}),
और उनका योग ही \emph{पथ पर कार्य}\index{कार्य (बल का)!पथ पर} है, $W_{AB}(\vect F) = \sum \scal{F}{\dd\ell}$ ---
जो अचर बल वाले सीधे पथ पर सिमटकर पिछले वर्ष का $F \times AB \times \cos\theta$ बन जाता है।
\end{definition}

\begin{omfigure}
\begin{tikzpicture}[font=\small]
  \draw[omExo!60, thick] plot [smooth, tension=0.8] coordinates {(0,0.3) (1.2,1.15) (2.6,1.3) (4.0,0.7) (5.3,1.05)};
  \draw[omMeth, thick] (0,0.3) -- (1.2,1.15) -- (2.6,1.3) -- (4.0,0.7) -- (5.3,1.05);
  \foreach \p in {(0,0.3),(1.2,1.15),(2.6,1.3),(4.0,0.7),(5.3,1.05)} \fill \p circle (1.3pt);
  \node[below] at (0,0.3) {$A$}; \node[below right] at (5.3,1.05) {$B$};
  \draw[-Stealth, omThm, very thick] (2.6,1.3) -- (3.44,0.94) node[below left=-2pt] {$\vect{\dd\ell}$};
  \draw[-Stealth, omDef, very thick] (2.6,1.3) -- ++(22:1.5) node[above] {$\vect F$};
  \draw (2.6,1.3) ++(-23:0.75) arc[start angle=-23, end angle=22, radius=0.75];
  \node at ($(2.6,1.3)+(0:1.0)$) {$\theta$};
\end{tikzpicture}

{\small मुड़ा हुआ पथ कई छोटे क़दमों की टूटी रेखा है: हर $\vect{\dd\ell}$ $F\,\dd\ell\cos\theta$ कमाता है;
और कार्य उनका योग है।}
\end{omfigure}

\begin{proposition}[भार का कार्य, किसी भी पथ पर]\label{prop:g12:work-and-energy:weight}
$A$ (ऊँचाई $z_A$) से $B$ (ऊँचाई $z_B$) तक \emph{किसी भी} पथ पर
द्रव्यमान $m$ का भार कार्य $W_{AB}(\vect P) = mg\,(z_A - z_B)$ करता है: हिसाब में गिरावट
आती है, रास्ता कभी नहीं।
\end{proposition}

\begin{proof}
हर छोटा क़दम $mg \times (\text{छोटी गिरावट})$ का योगदान देता है (पिछले वर्ष की सीधे-खंड वाली गणना); और
गिरावटें दूरबीन की तरह सिकुड़कर $z_A - z_B$ बन जाती हैं। उस समय मान ली गई वक्र वाली
स्थिति यहाँ चुका दी गई।
\end{proof}

\begin{definition}[संरक्षी और असंरक्षी बल]\label{def:g12:work-and-energy:conservative}
कोई बल \emph{संरक्षी}\index{संरक्षी बल} है यदि दो बिंदुओं के बीच उसका
कार्य हर पथ पर एक ही हो (और हर आने-जाने के चक्कर में शून्य), वरना वह
\emph{असंरक्षी}\index{असंरक्षी बल} है। भार और स्प्रिंग का बल
(अगला अनुभाग): संरक्षी। सरकन \omterm{def:g10:inertia:inventory}{घर्षण} नहीं: गति के विपरीत होने के कारण वह
लंबाई $L$ के पथ पर $W = -fL$ करता है --- यानी हर \omterm{def:g10:orders-of-magnitude:unit}{मीटर} की चुंगी, जो कभी
नहीं लौटती।
\end{definition}

\begin{example}[कुटिया तक दो पगडंडियाँ]\label{ex:g12:work-and-energy:trails}
\qty{75}{kg} का पदयात्री \qty{300}{m} चढ़ता है, या तो \qty{800}{m} की खड़ी पगडंडी से या \qty{2.4}{km} के
घुमावदार मोड़ों से। भार दोनों पर $-75 \times 9.81 \times 300 \approx \qty{-2.2e5}{J}$ वसूलता है; और \qty{30}{N} का
घर्षण-बल एक पर \qty{-2.4e4}{J}, दूसरे पर \qty{-7.2e4}{J}। उतरते समय भार
लौटा देता है; पर \omterm{def:g10:inertia:inventory}{घर्षण} का बिल दोनों बार ऊष्मा बनकर चला जाता है।
\end{example}

\section{स्थितिज ऊर्जा}

\begin{definition}[स्थितिज ऊर्जा]\label{def:g12:work-and-energy:ep}
हर \omterm{def:g12:work-and-energy:conservative}{संरक्षी बल} के साथ एक \emph{स्थितिज ऊर्जा}\index{स्थितिज ऊर्जा!संरक्षी
कार्य से} $E_p$ जोड़ दी जाती है, यानी केवल स्थिति पर निर्भर एक संख्या, जिसकी
परिभाषा $W_{AB} = E_p(A) - E_p(B)$ है: दिया गया कार्य ही ख़र्च हुई स्थितिज ऊर्जा
है। किसी चुने हुए संदर्भ बिंदु पर $E_p = 0$ तय कर देने से वह हर जगह तय हो जाती
है; और यह चुनाव सारे मानों को एक अचर जितना खिसका देता है, जो हर अंतर में से कट
जाता है। भार के लिए इससे पिछले वर्ष वाला $E_p = mgz$ फिर मिल जाता है, अब हर
पथ पर लागू। \omterm{def:g10:inertia:inventory}{घर्षण} का कोई $E_p$ हो ही नहीं सकता: पथ पर निर्भर बिल दो
सिरों की संख्याओं का अंतर नहीं होता --- संचित और वापस पाने लायक़ केवल
\omterm{def:g12:work-and-energy:conservative}{संरक्षी बल} रखते हैं, और नाम भी वही बात संरक्षित रखता है।
\end{definition}

\begin{proposition}[प्रत्यास्थ स्थितिज ऊर्जा]\label{prop:g12:work-and-energy:elastic}
\omterm{def:g12:oscillators-and-time:spring}{दृढ़ता} $k$ (\unit{N/m}) की स्प्रिंग $x$ जितनी खिंचने या दबने पर
बल $F = kx$ से वापस खींचती है (\cref{ch:g12:oscillators-and-time}); और \omterm{def:g10:relative-motion:frame}{विराम} पर लौटते समय वह जो
कार्य देती है --- यानी उसकी \emph{प्रत्यास्थ स्थितिज
ऊर्जा}\index{स्थितिज ऊर्जा!प्रत्यास्थ} --- $E_p = \tfrac12\, k x^2$ है।
\end{proposition}

\begin{proof}
\omterm{def:g12:oscillators-and-time:pendulum}{प्रत्यानयन बल} गति के अनुदिश लगता है, इसलिए $\sum F\,\dd\ell$ \emph{$F$--$x$
आलेख के नीचे का क्षेत्रफल} है: आधार $x$ और ऊँचाई $kx$ का त्रिभुज,
क्षेत्रफल $\tfrac12 x \times kx$। और हर खिंचाव पर बल जाते समय भी वही है और लौटते समय
भी: यानी कार्य सिरों से तय, स्प्रिंग का बल \omterm{def:g12:work-and-energy:conservative}{संरक्षी}, और
$E_p$ ठीक-ठीक परिभाषित।
\end{proof}

\begin{omfigure}
\begin{tikzpicture}
\begin{axis}[omaxis, width=8.4cm, height=4.8cm, xmin=0, xmax=0.095,
    ymin=0, ymax=21,
    xlabel={खिंचाव $x$ (m)}, ylabel={बल $F$ (N)}]
  \addplot[draw=none, fill=omDef!18, domain=0:0.06, samples=2] {200*x} \closedcycle;
  \addplot[omThm, thick, domain=0:0.09, samples=2] {200*x};
  \draw[dashed, omExo] (axis cs:0,12) -- (axis cs:0.06,12) -- (axis cs:0.06,0);
  \node[font=\small, anchor=west] at (axis cs:0.061,6.5) {$kx$};
  \node[font=\small] at (axis cs:0.044,3.2) {$\tfrac12 kx^2$};
  \node[font=\small, anchor=south, omThm] at (axis cs:0.082,17) {$F = kx$};
\end{axis}
\end{tikzpicture}

{\small स्प्रिंग का बल खिंचाव के साथ बढ़ता है; और संचित ऊर्जा
रेखा के नीचे का क्षेत्रफल है: एक त्रिभुज, $\tfrac12 \times x \times kx$।}
\end{omfigure}

\section{ऊर्जा की दो प्रमेय}

\begin{theorem}[गतिज-ऊर्जा प्रमेय]\label{thm:g12:work-and-energy:ketheorem}\index{गतिज-ऊर्जा प्रमेय}
द्रव्यमान $m$ के पिंड की $A$ से $B$ तक की किसी भी गति के लिए ---
मुड़ी हुई, फंदेदार, त्रिविमीय ---
\[
\Delta E_k = E_k(B) - E_k(A) = \sum W_{AB}(\vect F),
\]
जहाँ योग पिंड पर लगते सभी बलों पर है।
\end{theorem}

\begin{proof}
चूँकि $v^2 = \scal{v}{v}$, इस वर्ष के गणित का गुणनफल-नियम $v^2$ का अवकलज $2\,\scal{a}{v}$ दे देता
है, इसलिए \omterm{thm:g12:newtons-laws:second}{न्यूटन के दूसरे नियम} (\cref{ch:g12:newtons-laws}) से $\dd E_k/\dd t = m\,\scal{a}{v} = \scal{F_{\text{कुल}}}{v}$। छोटे $\dd t$ में
पिंड $\vect{\dd\ell} = \vect v\,\dd t$ चलता है: $E_k$ \cref{def:g12:work-and-energy:path} का छोटा कार्य कमा लेता है; और
क़दम जोड़ने से कार्य जुड़ जाते हैं --- यानी पिछले वर्ष मान लिया गया
कथन, पूरा-पूरा कमाया हुआ।
\end{proof}

\begin{proposition}[तात्क्षणिक शक्ति]\label{prop:g12:work-and-energy:power}
हर क्षण वेग $\vect v$ वाले पिंड पर लगता बल $\vect F$ शक्ति
$P = \scal{F}{v} = Fv\cos\theta$ देता है, और $\dd E_k/\dd t$ सभी बलों की शक्तियों का योग है:
यानी पिछले वर्ष का अचर-वेग वाला सूत्र, किसी भी गति के हर क्षण पर ठीक-ठीक।
\end{proposition}

\begin{proof}
पिछली उपपत्ति से पढ़ लीजिए: हर बल $\dd t$ के दौरान $E_k$ में
$\scal{F}{\dd\ell} = \scal{F}{v}\,\dd t$ डालता है।
\end{proof}

\begin{theorem}[यांत्रिक-ऊर्जा प्रमेय]\label{thm:g12:work-and-energy:emtheorem}\index{यांत्रिक-ऊर्जा प्रमेय}
मान लीजिए $E_m = E_k + E_p$, जहाँ $E_p$ कार्य कर रहे सभी \omterm{def:g12:work-and-energy:conservative}{संरक्षी}
बलों की स्थितिज ऊर्जाएँ समेट लेता है। तब, किसी भी पथ पर,
\[
\Delta E_m = W_{AB}(\text{असंरक्षी बल}),
\]
और जब भी \omterm{def:g12:work-and-energy:conservative}{असंरक्षी बल} कोई कार्य न करें, $E_m$ संरक्षित रहता
है।
\end{theorem}

\begin{proof}
गतिज-ऊर्जा प्रमेय को बाँट दीजिए: $\Delta E_k = W_{\text{संरक्षी}} +
W_{\text{असंरक्षी}} = -\Delta E_p + W_{\text{असंरक्षी}}$; और $\Delta E_p$ को दूसरी ओर ले
जाइए।
\end{proof}

\begin{example}[लोलक, इस बार ईमानदारी से]\label{ex:g12:work-and-energy:pendulum}
लंबाई $L = \qty{2.5}{m}$ के तार पर लटका लोलक का गोलक \omterm{def:g10:universal-gravitation:weight}{ऊर्ध्वाधर} से $\qty{45}{\degree}$ पर \omterm{def:g10:relative-motion:frame}{विराम} में
छोड़ा जाता है। गति के लंबवत होने के कारण \omterm{def:g10:inertia:inventory}{तनाव} शक्तिहीन है; और
भार \omterm{def:g12:work-and-energy:conservative}{संरक्षी} है: इसलिए $E_m$ \emph{मुड़े हुए चाप पर}
संरक्षित रहता है, जो पिछले वर्ष केवल जाँचा जा सकता था। गोलक $h = L(1 - \cos\qty{45}{\degree}) =
\qty{0.73}{m}$ ऊपर से
चलता है, इसलिए तल पर $v = \sqrt{2gh} \approx
\qty{3.8}{m/s}$।
\end{example}

\begin{method}[किसी भी गति का ऊर्जा-लेखा]\label{met:g12:work-and-energy:audit}
\begin{enumerate}
  \item निकाय चुनिए और उस पर लगते बल गिनाइए।
  \item छाँटिए: गति के लंबवत $\to$ कोई कार्य नहीं; \omterm{def:g12:work-and-energy:conservative}{संरक्षी}
        $\to$ $E_m$ में एक $E_p$; और बाक़ी $\to$ $W_{\text{असंरक्षी}}$ (सरकन
        \omterm{def:g10:inertia:inventory}{घर्षण}: $-fL$, $L$ = पथ-लंबाई)।
  \item आँकड़े वाले बिंदु और प्रश्न वाले बिंदु के बीच $\Delta E_m = W_{\text{असंरक्षी}}$ लिखिए; और हल
        कीजिए। ऊर्जा ``कितनी तेज़, कितनी ऊँची'' का उत्तर देती है;
        केवल समय और त्वरण के लिए न्यूटन के नियम चाहिए।
\end{enumerate}
\end{method}

\section{ऊर्जा आरेख}

\begin{definition}[ऊर्जा आरेख, पलटाव बिंदु]\label{def:g12:work-and-energy:diagram}
जिस गति में एक ही अक्ष $x$ हो और $E_m$ संरक्षित रहे, उसके लिए
\emph{ऊर्जा आरेख}\index{ऊर्जा आरेख} वक्र $E_p(x)$ और क्षैतिज रेखा $E_m$ को
आलेखित करता है। चूँकि $E_k = E_m - E_p \geq 0$, गति केवल वहीं तक सीमित रहती है जहाँ वक्र रेखा
से नीचे हो; और दोनों के बीच का अंतर सीधे $E_k$ पढ़ लिया जाता है। जहाँ वक्र
रेखा से मिल जाए, यानी \emph{पलटाव बिंदु}\index{पलटाव बिंदु} पर, $v = 0$ होता
है और गति उलट जाती है।
\end{definition}

\begin{omfigure}
\begin{tikzpicture}
\begin{axis}[omaxis, width=9.0cm, height=5.4cm,
    xlabel={$x$}, ylabel={ऊर्जा}, xtick=\empty, ytick=\empty,
    xmin=-2.25, xmax=2.25, ymin=0, ymax=2.45]
  \addplot[omDef, thick, domain=-2.05:2.05, samples=80] {0.25*x^4 - x^2 + 1.5};
  \addplot[omThm, dashed, thick, domain=-2.25:2.25, samples=2] {1.0};
  \addplot[omProp, dashed, domain=-2.25:2.25, samples=2] {1.7};
  \node[font=\small, omThm, anchor=west] at (axis cs:-2.2,0.86) {$E_m$};
  \node[font=\small, omProp, anchor=west] at (axis cs:-2.2,1.85) {$E_m'$};
  \node[font=\small, omDef, anchor=west] at (axis cs:1.55,2.1) {$E_p(x)$};
  \fill (axis cs:0.765,1.0) circle (1.6pt) node[font=\small, below=1pt] {$x_1$};
  \fill (axis cs:1.848,1.0) circle (1.6pt) node[font=\small, below right=-1pt] {$x_2$};
  \fill[omMeth] (axis cs:1.414,0.5) circle (2.2pt) node[font=\small, below] {स्थायी};
  \fill[omMeth] (axis cs:0,1.5) circle (2.2pt) node[font=\small, above] {अस्थायी};
\end{axis}
\end{tikzpicture}

{\small एक \omterm{def:g12:work-and-energy:diagram}{ऊर्जा आरेख}, बिना कुछ हल किए पढ़ा हुआ: ऊर्जा
$E_m$ के साथ पिंड $x_1$ और $x_2$ के बीच दोलन करता है, और वहाँ सबसे तेज़
जहाँ अंतर सबसे बड़ा हो; और $E_m'$ के साथ वह पहाड़ी को रेंगकर पार कर जाता है और
\omterm{rem:g12:work-and-energy:escape}{भाग निकलता} है।}
\end{omfigure}

\begin{proposition}[भू-दृश्य से बल; साम्यावस्थाएँ]\label{prop:g12:work-and-energy:equilibria}
$x$ के अनुदिश \omterm{def:g12:work-and-energy:conservative}{संरक्षी बल} $E_p$ की ढाल का ऋणात्मक है: $F = -\,\dd E_p/\dd x$
--- यानी वह आरेख पर \emph{नीचे की ओर} संकेत करता है और वहाँ लुप्त हो जाता है
जहाँ स्पर्शी क्षैतिज हो: यानी साम्यावस्था पर। $E_p$ का निम्निष्ठ
\emph{स्थायी}\index{स्थायी साम्यावस्था} है --- हटाया गया पिंड वापस धकेल दिया
जाता है, जैसे घाटी में गेंद; और उच्चिष्ठ \emph{अस्थायी}\index{अस्थायी
साम्यावस्था} --- वह दूर धकेल दिया जाता है, जैसे पहाड़ी की चोटी पर गेंद।
\end{proposition}

\begin{proof}
छोटे क़दम $\dd x$ पर बल $F\,\dd x = E_p(x) -
E_p(x + \dd x) = -\dd E_p$ करता है; $\dd x$ से भाग दे
दीजिए। घाटी के दोनों ओर नीचे की दिशा भीतर की ओर संकेत करती है; और पहाड़ी की
चोटी के चारों ओर बाहर की ओर।
\end{proof}

\begin{remark}[दोलन और पलायन]\label{rem:g12:work-and-energy:escape}
\cref{ch:g12:oscillators-and-time} के छोटे दोलन $E_p$ की घाटियों के तलों में रहते हैं। और पृथ्वी से
दूर गुरुत्व कमज़ोर पड़ता जाता है: $E_p(r)$ का वक्र लगातार धीमे-धीमे एक परिमित छत
की ओर चढ़ता है (सिकुड़ते बल-आलेख के नीचे का क्षेत्रफल परिमित है)। जिस राकेट का
$E_m$ उस छत तक पहुँच जाए उसे कोई \omterm{def:g12:work-and-energy:diagram}{पलटाव बिंदु} मिलता ही नहीं: वह
\emph{भाग निकलता}\index{पलायन चाल} है --- ज़मीन से लगभग \qty{11.2}{km/s} पर (\cref{exo:g12:work-and-energy:13});
और उससे नीचे वह बँधा रहता है। छत ख़ुद वर्ष 1 के खंड में निकाली गई है।
\end{remark}

\begin{example}[फंदे वाला चक्कर]\label{ex:g12:work-and-energy:loop}
ऊँचाई $H$ से छोड़ा गया कंचा त्रिज्या $R$ के \omterm{def:g10:universal-gravitation:weight}{ऊर्ध्वाधर} फंदे में लुढ़क
जाता है। वृत्तीय गतिकी (\cref{ch:g12:newtons-laws}) माँगती है कि शिखर पर $v_{\text{शिखर}}^2 \geq gR$ --- यानी वही
आँकड़ा जो पिछले वर्ष की रोलर-कोस्टर समस्या ने उधार लिया था, अब हमारा अपना।
$H$ पर \omterm{def:g10:relative-motion:frame}{विराम} से शिखर (ऊँचाई $2R$) तक संरक्षण $v_{\text{शिखर}}^2 = 2g(H - 2R)$ देता है, इसलिए
$H \geq \tfrac52 R$: यानी फंदे के शिखर से आधी त्रिज्या ऊपर, द्रव्यमान चाहे जो हो --- और
असली, रगड़ती दुनिया में उससे थोड़ा और ऊपर।
\end{example}

\begin{omfigure}
\begin{tikzpicture}[font=\small, scale=0.9]
  \draw[thick] (0,3.0) .. controls (0.9,1.1) and (1.6,0) .. (3.1,0) -- (6.9,0);
  \draw[thick] (5.5,1.2) circle (1.2);
  \fill[omDef] (0,3.0) circle (2.2pt);
  \draw[dashed, omExo] (4.0,2.4) -- (7.6,2.4) node[right] {$2R$};
  \draw[dashed, omThm] (-0.4,3.0) -- (7.6,3.0) node[right] {$H_{\min} = \tfrac52 R$};
  \draw[Stealth-Stealth, omThm] (7.35,2.4) -- (7.35,3.0);
  \node[right, omThm] at (7.35,2.7) {$\tfrac{R}{2}$};
  \draw[-Stealth, omMeth, thick] (5.5,2.4) -- (4.5,2.4) node[above right=0pt and 2pt] {$v_{\text{शिखर}}$};
  \fill (5.5,2.4) circle (1.5pt);
\end{tikzpicture}

{\small फंदे वाला चक्कर: छोड़ने का बिंदु शिखर से आधी त्रिज्या ऊपर बैठता है,
इसलिए बची हुई $E_k$ कंचे को दबा हुआ बनाए रखती है ($v_{\text{शिखर}}^2 = gR$)।}
\end{omfigure}

\section{अभ्यास}

\begin{exercise}[$\star$]\label{exo:g12:work-and-energy:1}
$k = \qty{200}{N/m}$ की \omterm{def:g12:oscillators-and-time:spring}{दृढ़ता} वाली स्प्रिंग \qty{5.0}{cm} जितनी खींची जाती है। संचित
ऊर्जा निकालिए; खिंचाव दुगुना करने पर वह क्या हो जाती है? और
\emph{दूसरे} \qty{5.0}{cm} का दाम कितना पड़ता है?
\end{exercise}

\begin{exercise}[$\star$]\label{exo:g12:work-and-energy:2}
\qty{0.90}{kg} का ड्रोन $z = \qty{12}{m}$ पर बनी छत से $z = \qty{30}{m}$ पर बनी बालकनी तक घूमता-फिरता
रास्ता उड़ता है। उसके भार का कार्य? सीधे चढ़ाव पर? आने-जाने के
चक्कर पर?
\end{exercise}

\begin{exercise}[$\star$]\label{exo:g12:work-and-energy:3}
\omterm{def:g12:work-and-energy:conservative}{संरक्षी}, \omterm{def:g12:work-and-energy:conservative}{असंरक्षी}, या कार्यहीन? पथों के आधार पर कारण दीजिए:
(क) भार; (ख) स्प्रिंग का बल; (ग) सरकन \omterm{def:g10:inertia:inventory}{घर्षण};
(घ) सरकते डिब्बे पर \omterm{def:g12:newtons-laws:friction}{अभिलंब बल}; (ङ) वायु-कर्षण।
\end{exercise}

\begin{exercise}[$\star$]\label{exo:g12:work-and-energy:4}
लंबाई \qty{2.0}{m} का लोलक \omterm{def:g10:universal-gravitation:weight}{ऊर्ध्वाधर} से $\qty{60}{\degree}$ पर \omterm{def:g10:relative-motion:frame}{विराम} में छोड़ा जाता है। सबसे
निचले बिंदु पर उसकी चाल निकालिए। तार का \omterm{def:g10:inertia:inventory}{तनाव} इस तराज़ू में कभी क्यों
नहीं आता?
\end{exercise}

\begin{exercise}[$\star$]\label{exo:g12:work-and-energy:5}
कोई साइकिल-सवार समतल पर \qty{9.0}{m/s} से चलती है और \qty{250}{W} देती है: वह कुल कितने
\omterm{def:g11:work-of-force:sign}{प्रतिरोधी} बल से लड़ रही है? उसी शक्ति के साथ
\qty{5.0}{m/s} चढ़ते समय: वह किस बल पर विजय पाती है?
\end{exercise}

\begin{exercise}[$\star\star$]\label{exo:g12:work-and-energy:6}
किसी खिलौना-प्रक्षेपक की स्प्रिंग ($k = \qty{450}{N/m}$) \qty{20}{g} के तीर के पीछे \qty{8.0}{cm}
जितनी दबाई जाती है। संचित ऊर्जा, छोड़ने की चाल, और सीधे ऊपर दागने पर
पहुँची ऊँचाई निकालिए (हवा नहीं)।
\end{exercise}

\begin{exercise}[$\star\star$]\label{exo:g12:work-and-energy:7}
रबर बैंड खींचते हुए कोई विद्यार्थी $x = 0$, $2.0$, $4.0$, $6.0$, \qty{8.0}{cm}
पर $F = 0$, $3.0$, $7.0$, $12.0$, \qty{18.0}{N} नापता है। समलंब-क्षेत्रफलों से
संचित ऊर्जा का अनुमान लगाइए; अंत में $k = F/x$ लेकर उसकी तुलना $\tfrac12 kx^2$
से कीजिए। दोनों अलग क्यों हैं?
\end{exercise}

\begin{exercise}[$\star\star$]\label{exo:g12:work-and-energy:8}
\qty{55}{kg} की स्लेज \omterm{def:g10:relative-motion:frame}{विराम} से चलकर \qty{25}{m} लंबे मुड़े हुए रास्ते पर \qty{8.0}{m} गिरती
है और \qty{9.0}{m/s} पर पहुँचती है। कितनी \omterm{def:g11:mechanical-energy:em}{यांत्रिक ऊर्जा} गँवाई गई? औसत
घर्षण-बल? रास्ते का ठीक आकार बेमतलब क्यों है?
\end{exercise}

\begin{exercise}[$\star\star$]\label{exo:g12:work-and-energy:9}
कंचे की पटरी में त्रिज्या \qty{20}{cm} का फंदा है। न्यूनतम छोड़-ऊँचाई निकालिए
(\omterm{def:g10:relative-motion:frame}{विराम} से, कोई \omterm{def:g10:inertia:inventory}{घर्षण} नहीं)। \qty{60}{cm} से छोड़ने पर: फंदे के शिखर पर चाल,
और अनुपात $v_{\text{शिखर}}^2/(gR)$? असली पटरी को उससे भी ऊँचाई से क्यों चलाना पड़ता है?
\end{exercise}

\begin{exercise}[$\star\star$]\label{exo:g12:work-and-energy:10}
\qty{100}{g} का मनका बिना \omterm{def:g10:inertia:inventory}{घर्षण} किसी अक्ष पर सरकता है, जहाँ $E_p(x) = 8.0\,x^2$
(\omterm{def:g11:work-of-force:work}{जूल}, $x$ \omterm{def:g10:orders-of-magnitude:unit}{मीटर} में) और $E_m = \qty{2.0}{J}$। \omterm{def:g12:work-and-energy:diagram}{पलटाव बिंदु}?
अधिकतम चाल? यह किस तरह की गति है?
\end{exercise}

\begin{exercise}[$\star\star$]\label{exo:g12:work-and-energy:11}
\qty{2.0}{m/s} पर चलती \qty{1200}{kg} की कार \omterm{def:g12:oscillators-and-time:spring}{दृढ़ता} $k = \qty{6.0e5}{N/m}$ वाली बंपर-स्प्रिंग से
रोकी जाती है। अधिकतम \omterm{def:g12:sound-acoustics:sound-wave}{संपीडन} निकालिए। \qty{4.0}{m/s} पर: ऊर्जा
चौगुनी होने पर भी वह चार गुना क्यों नहीं है?
\end{exercise}

\begin{exercise}[$\star\star\star$]\label{exo:g12:work-and-energy:12}
\qty{0.50}{kg} की ठेली ऐसी पटरी पर लुढ़कती है जिसके $E_p(x)$ में \qty{6.0}{J} की पहाड़ी
($x = 0$) और \qty{2.0}{J} की घाटी ($x = \qty{1.5}{m}$) है। (क) बाईं ओर से $E_m = \qty{5.0}{J}$ के साथ आती
है: क्या वह पार कर जाती है? नहीं तो क्या होता है? (ख) $E_m = \qty{7.0}{J}$ के साथ: पहाड़ी
और घाटी पर चालें? (ग) साम्यावस्थाएँ?
\end{exercise}

\begin{exercise}[$\star\star\star$]\label{exo:g12:work-and-energy:13}
यह मानते हुए कि द्रव्यमान $m$ को ज़मीन से मनमाने ढंग से दूर तक ढोने में
$mgR_E$, $R_E = \qty{6.37e6}{m}$ लगता है (यानी कमज़ोर पड़ते गुरुत्व-आलेख के नीचे का परिमित
क्षेत्रफल), पृथ्वी से \emph{\omterm{rem:g12:work-and-energy:escape}{पलायन चाल}} निकालिए। वह यान और कंकड़, दोनों
के लिए एक ही क्यों है? और उससे ज़रा नीचे के प्रक्षेपण का क्या होगा?
\end{exercise}

\begin{exercise}[$\star\star\star$]\label{exo:g12:work-and-energy:14}
\qty{250}{g} का फिसलक स्प्रिंग ($k = \qty{40}{N/m}$) पर \qty{6.0}{cm} (\cref{ch:g12:oscillators-and-time}) के आयाम से
दोलन करता है। कुल ऊर्जा? अधिकतम चाल? वे स्थितियाँ जहाँ $E_k = E_p$?
\end{exercise}

\begin{exercise}[$\star\star\star$]\label{exo:g12:work-and-energy:15}
कोई बाँस-कूद खिलाड़ी \qty{9.5}{m/s} से दौड़ता है। यदि बाँस उसकी पूरी \omterm{def:g11:mechanical-energy:kinetic}{गतिज
ऊर्जा} बदल दे तो उसका \omterm{def:g11:forces-and-motion:com}{द्रव्यमान केंद्र} कितनी ऊँचाई कमाएगा, और छड़ कितनी पार
होगी (\omterm{def:g11:forces-and-motion:com}{द्रव्यमान केंद्र} \qty{1.0}{m} ऊपर से शुरू होकर), इसका अनुमान लगाइए। रिकॉर्ड
लगभग \qty{6.3}{m} है: यह अनुमान क्या पकड़ लेता है, और क्या छोड़ देता है?
\end{exercise}

\section{समस्या: स्की-कूद}

\begin{problem}[{सप्ताहांत समस्या --- स्की-कूद: चाल-मापी बंदूक़ से जाँची गई
ढलान, मुक्त पतन की उड़ान, खड़ी लैंडिंग की कोमल ज्यामिति, और वह मीनार जो
घर्षण का बिल आने पर डिज़ाइनर को ख़रीदनी पड़ती है}]
\label{pb:g12:work-and-energy:1}
कोई ओलंपिक स्की-कूद खिलाड़ी (स्कियों समेत $m = \qty{70}{kg}$) उछाल-मेज़ से $H = \qty{45}{m}$ ऊपर
बने फाटक से \omterm{def:g10:relative-motion:frame}{विराम} में चलता है, $L = \qty{90}{m}$ की मुड़ी हुई ढलान पर सरकता है, और मेज़
से क्षैतिज दिशा में निकल जाता है; वहाँ लगी चाल-मापी बंदूक़ $v_0 = \qty{26}{m/s}$ पढ़ती है।
उतरने का बिंदु $h = \qty{40}{m}$ नीचे, $\qty{35}{\degree}$ की ढलान पर है। भाग II और III में
वायु-प्रतिरोध छोड़ दिया गया है।

\medskip\noindent\textbf{भाग I --- ढलान।}
\begin{enumerate}
  \item घर्षणहीन उछाल-चाल का पूर्वानुमान लगाइए; ढलान का वक्र बेमतलब क्यों है?
  \item संदर्भ मेज़ पर रखकर फाटक पर और उछाल पर $E_m$ निकालिए।
  \item $\Delta E_m$ निकालिए; कौन-सी प्रमेय दोषी बताती है, और दोषी कौन है?
  \item पूरी ढलान पर औसत \omterm{def:g11:work-of-force:sign}{प्रतिरोधी} बल (बर्फ़ और हवा) निकालिए।
  \item $E_m$ का कितना अंश खो जाता है? स्कियों पर मोम क्यों लगाया जाता है
        और खिलाड़ी क्यों दुबक जाता है?
\end{enumerate}

\medskip\noindent\textbf{भाग II --- उड़ान।}
\begin{enumerate}[resume]
  \item उड़ान में केवल भार लगता है: यह \cref{ch:g12:projectile-motion} की कौन-सी गति है?
        उतरते समय वेग के घटक?
  \item उतरने की चाल दो बार निकालिए: घटकों से, और $E_m$ के संरक्षण से।
        \unit{km/h} में?
  \item उड़ान का समय और उड़ी गई क्षैतिज दूरी निकालिए।
  \item असली खिलाड़ी हवा पर पंख की तरह सवार होकर कहीं दूर तक उड़ते हैं: क्या
        हवा उतरने की चाल \emph{बढ़ा} सकती है? यांत्रिक-ऊर्जा प्रमेय
        से तर्क दीजिए।
\end{enumerate}

\medskip\noindent\textbf{भाग III --- उतरना।}
\begin{enumerate}[resume]
  \item उतरते वेग का क्षैतिज से नीचे का कोण निकालिए।
  \item उसे ढलान के समांतर और लंबवत घटकों में बाँटिए।
  \item टाँगें केवल लंबवत हिस्सा सोखती हैं: उस \omterm{def:g11:mechanical-energy:kinetic}{गतिज ऊर्जा} को निकालिए;
        कुल से तुलना कीजिए।
  \item समतल ज़मीन पूरी की पूरी सोख लेती: हर लैंडिंग किस \omterm{def:g10:universal-gravitation:weight}{ऊर्ध्वाधर} गिरावट के
        तुल्य है ($h_{\text{तुल्य}} = v_\perp^2/2g$, और $v^2/2g$)?
  \item एक वाक्य में: उतरने की पहाड़ियाँ खड़ी और मुड़ी हुई क्यों होती हैं?
\end{enumerate}

\medskip\noindent\textbf{भाग IV --- डिज़ाइनर की मीनार।}
किसी छोटी पहाड़ी को उछाल पर केवल $v_0' = \qty{24}{m/s}$ चाहिए; उसकी सीधी ढलान $\qty{35}{\degree}$ पर
उतरती है, इसलिए फाटक-ऊँचाई $H'$ का अर्थ पटरी-लंबाई $H'/\sin\qty{35}{\degree}$ है; और कर्षण
वही \qty{80}{N}।
\begin{enumerate}[resume]
  \item कोई घर्षणहीन डिज़ाइनर जितनी फाटक-ऊँचाई बनाता, वह निकालिए।
  \item $H'$ के लिए फाटक से मेज़ तक यांत्रिक-ऊर्जा प्रमेय लिखिए।
  \item $H'$ के लिए हल कीजिए।
  \item \omterm{def:g10:inertia:inventory}{घर्षण} ने मीनार कितने प्रतिशत ऊँची कर दी?
  \item अब द्रव्यमान कटता नहीं: \qty{60}{kg} के खिलाड़ी के लिए $H'$ फिर से
        निकालिए। ऊँची मीनार किसे चाहिए, और क्यों?
  \item अंतिम चोट: घोषित की गई मीनार, और उसमें \omterm{def:g10:inertia:inventory}{घर्षण} का हिस्सा।
\end{enumerate}
\end{problem}
